% Appendix B: Source Code

\section{Repository Structure}

The complete source code is available at:

\begin{verbatim}
https://github.com/marufdevops/multiagent-cooperation-research
\end{verbatim}

\subsection{Directory Organization}

\begin{verbatim}
Dissertation/
├── src/
│   ├── agents/
│   │   ├── harvester_agent.py
│   │   └── fruit.py
│   ├── models/
│   │   └── harvest_model.py
│   ├── visualization/
│   │   └── solara_viz.py
│   ├── analysis/
│   │   └── statistical_analysis.py
│   └── utils/
│       └── metrics.py
├── scripts/
│   ├── run_sandbox.py
│   ├── run_smoke_test.py
│   ├── run_full_experiment.py
│   └── run_visualization.py
├── config/
│   └── experiment_config.yaml
├── tests/
│   └── test_model.py
├── dissertation/
│   ├── main.tex
│   ├── chapters/
│   └── figures/
├── requirements.txt
└── README.md
\end{verbatim}

\section{Key Code Excerpts}

\subsection{Harvester Agent Implementation}

\begin{lstlisting}[language=Python, caption=HarvesterAgent class (excerpt)]
class HarvesterAgent(Agent):
    """Agent that moves, harvests fruit, and communicates."""
    
    def __init__(self, model):
        super().__init__(model)
        self.harvested = 0
        self.messages_sent = 0
        self.messages_received = 0
        self.current_target = None
        
    def step(self):
        """Execute one step: move, harvest, communicate."""
        self.move()
        self.harvest()
        self.communicate()
        
    def communicate(self):
        """Share information with agents within comm_range."""
        if self.model.comm_range == 0:
            return
        
        neighbors = self.model.grid.get_neighbors(
            self.pos,
            moore=True,
            radius=self.model.comm_range
        )
        
        nearest_fruit = self._find_nearest_fruit()
        if nearest_fruit:
            for neighbor in neighbors:
                if isinstance(neighbor, HarvesterAgent):
                    neighbor.receive_message(nearest_fruit)
                    self.messages_sent += 1
\end{lstlisting}

\subsection{Model Implementation}

\begin{lstlisting}[language=Python, caption=HarvestModel class (excerpt)]
class HarvestModel(Model):
    """Multi-agent fruit harvesting simulation."""
    
    def __init__(self, width=50, height=50, num_agents=10,
                 fruit_density=0.2, comm_range=2,
                 dynamics='Static', regen_prob=0.0, seed=None):
        super().__init__(seed=seed)
        
        self.width = width
        self.height = height
        self.num_agents = num_agents
        self.comm_range = comm_range
        self.dynamics = dynamics
        self.regen_prob = regen_prob
        
        self.grid = MultiGrid(width, height, torus=False)
        self._place_fruits()
        self._place_agents()
        
        self.datacollector = DataCollector(
            model_reporters={
                'total_yield': self._get_total_yield,
                'messages_this_step': self._get_messages_this_step,
                # ... additional metrics
            }
        )
        
    def step(self):
        """Execute one model step using Mesa 3.0 convention."""
        self.agents.shuffle_do('step')
        self.datacollector.collect(self)
\end{lstlisting}

\subsection{Experimental Runner}

\begin{lstlisting}[language=Python, caption=Factorial experiment runner (excerpt)]
def run_factorial_experiment(param_grid, output_dir):
    """Run full factorial experiment."""
    configs = generate_configurations(param_grid)
    results = []
    
    for config in tqdm(configs):
        model = HarvestModel(**config)
        model.run_model(steps=500)
        
        summary = extract_summary_metrics(model)
        results.append(summary)
    
    df = pd.DataFrame(results)
    df.to_csv(f'{output_dir}/results.csv', index=False)
    return df
\end{lstlisting}

\subsection{Statistical Analysis}

\begin{lstlisting}[language=Python, caption=ANOVA analysis (excerpt)]
import statsmodels.api as sm
from statsmodels.formula.api import ols

def run_anova(df, outcome='final_yield'):
    """Perform multi-way ANOVA."""
    formula = (f'{outcome} ~ C(comm_range) * C(num_agents) * '
               f'C(fruit_density) * C(dynamics)')
    
    model = ols(formula, data=df).fit()
    anova_table = sm.stats.anova_lm(model, typ=2)
    
    return anova_table
\end{lstlisting}

\section{Running the Code}

\subsection{Installation}

\begin{verbatim}
# Clone repository
git clone https://github.com/marufdevops/
    multiagent-cooperation-research.git
cd multiagent-cooperation-research

# Install dependencies
pip install -r requirements.txt
\end{verbatim}

\subsection{Running Experiments}

\begin{verbatim}
# Run sandbox model
python scripts/run_sandbox.py

# Run smoke test
python scripts/run_smoke_test.py

# Run full experiment
python scripts/run_full_experiment.py --config config/full.yaml

# Launch visualization
python scripts/run_visualization.py
# Then open browser to http://localhost:8765
\end{verbatim}

\subsection{Analysis}

\begin{verbatim}
# Run statistical analysis
python scripts/analyze_results.py --input results.csv

# Generate plots
python scripts/generate_plots.py --input results.csv
\end{verbatim}

\section{Software Versions}

\begin{table}[h]
\centering
\caption{Software Dependencies and Versions}
\begin{tabular}{ll}
\toprule
\textbf{Package} & \textbf{Version} \\
\midrule
Python & 3.10+ \\
Mesa & 3.0.0+ \\
Solara & 1.30.0+ \\
pandas & 2.0.0+ \\
NumPy & 1.24.0+ \\
SciPy & 1.10.0+ \\
statsmodels & 0.14.0+ \\
Matplotlib & 3.7.0+ \\
Seaborn & 0.12.0+ \\
\bottomrule
\end{tabular}
\end{table}

\section{License}

The code is released under the MIT License, allowing free use, modification, and distribution with attribution.

